\section{Các Công trình Liên quan}
\label{sec:tong-quan-tai-lieu}

\subsection{Phân loại pha giáng thủy}
\label{sec:lit-phan-loai-giang-thuy}

Hướng tiếp cận truyền thống nhất để phân biệt mưa và tuyết là áp một ngưỡng nhiệt độ
không khí hoặc nhiệt độ ướt (wet-bulb temperature) cố định hoặc theo vùng khí hậu
\cite{jennings2018}. Các nghiên cứu này cho thấy ngưỡng tối ưu thay đổi đáng kể theo vị
trí địa lý và mùa, và luôn tồn tại một dải nhiệt độ chồng lấn mà bản thân nhiệt độ không
đủ để phân biệt hai pha. Các công trình gần đây hơn \cite{webb2025} đặt vấn đề học máy
cho bài toán này, chỉ ra rằng mô hình học máy có thể khai thác thêm thông tin từ độ ẩm,
áp suất và gió để cải thiện độ chính xác trong vùng chồng lấn, nhưng phần lớn các nghiên
cứu này giả định dữ liệu đầu vào đầy đủ.

\textbf{Hệ quả cho đồ án:} ngưỡng nhiệt độ được dùng làm baseline tối thiểu
(Mục~\ref{sec:baseline}), còn các mô hình học máy/học sâu được kỳ vọng vượt baseline này
đặc biệt trong dải nhiệt độ chồng lấn và khi dữ liệu bị thiếu.

\subsection{Trích xuất đặc trưng chuỗi thời gian và mô hình cây}
\label{sec:lit-feature-tree}

Một hướng tiếp cận phổ biến cho phân loại chuỗi thời gian là trích xuất một tập đặc
trưng thống kê/thời gian-tần số từ mỗi cửa sổ, sau đó đưa vào một mô hình phân loại dạng
bảng. Các bộ công cụ tiêu biểu gồm tsfresh \cite{christ2018tsfresh}, trích xuất hàng trăm
đặc trưng tự động có kiểm định ý nghĩa thống kê, và catch22 \cite{lubba2019catch22}, một
tập 22 đặc trưng rút gọn được chọn để vừa có tính phân biệt cao vừa giảm chi phí tính
toán. Kết hợp với các mô hình gradient boosting như XGBoost \cite{chen2016xgboost}, các
phương pháp này thường đạt hiệu năng cạnh tranh với chi phí huấn luyện thấp hơn học sâu.

\textbf{Hệ quả cho đồ án:} nhánh Classic ML (Mục~\ref{sec:nhanh-classic-ml}) kết hợp đặc
trưng catch22/tsfresh với XGBoost, đồng thời khai thác khả năng nhận giá trị thiếu trực
tiếp của XGBoost như một baseline xử lý thiếu dữ liệu đơn giản.

\subsection{Họ ROCKET}
\label{sec:lit-rocket}

Họ ROCKET biến đổi chuỗi thời gian bằng các kernel tích chập và dùng đặc trưng tổng hợp để huấn luyện mô hình tuyến tính. MiniRocket \cite{dempster2021minirocket} sử dụng tập kernel phần lớn xác định, khai thác tỷ lệ phần tử dương (PPV) với các dilation và bias khác nhau. MultiRocket \cite{tan2022multirocket} mở rộng trên chuỗi gốc và sai phân bậc một, sử dụng nhiều phép pooling nhằm làm phong phú biểu diễn đặc trưng.

\textbf{Hệ quả cho đồ án:} cả hai cần đầu vào hữu hạn, nên nhánh ROCKET (Mục~\ref{sec:nhanh-rocket}) nội suy trong cửa sổ lịch sử và nối đặc trưng mask/$\Delta t$ để bổ sung thông tin vị trí thiếu. Thực nghiệm đánh giá các biến thể có ngân sách kernel giảm để phù hợp tài nguyên, không nhằm tái hiện đầy đủ benchmark gốc. Kết quả chỉ được so sánh trong giao thức Szeged của đồ án.

\subsection{RNN xử lý dữ liệu thiếu}
\label{sec:lit-rnn-missing}

GRU-D \cite{che2018grud} mở rộng GRU bằng cách đưa trực tiếp mặt nạ thiếu dữ liệu và thời
gian trôi qua từ lần quan trắc hợp lệ gần nhất ($\Delta t$) vào cổng decay của trạng thái
ẩn và của giá trị đầu vào, cho phép mô hình "quên" có kiểm soát khi dữ liệu thiếu kéo
dài. BRITS \cite{cao2018brits} đi theo hướng khác: xây dựng một mô hình imputation hai
chiều (forward/backward) dựa trên RNN, huấn luyện đồng thời với tác vụ phân loại để lan
truyền thông tin ngữ cảnh qua các đoạn thiếu.

\textbf{Hệ quả cho đồ án:} cả hai kiến trúc được triển khai trong nhánh Deep Learning
(Mục~\ref{sec:nhanh-deep-learning}) cùng với một mô hình GRU đơn giản có mask/$\Delta t$
làm đặc trưng bổ sung (GRU-Fill), nhằm đo đóng góp riêng của cơ chế decay có kiểm soát.

\subsection{Transformer và mô hình nền tảng}
\label{sec:lit-transformer}

SAITS \cite{du2023saits} dùng cơ chế self-attention với một mục tiêu kép (joint-optimization
giữa imputation và dự đoán) để tránh phải xây hai pha huấn luyện riêng biệt như BRITS.
Các mô hình nền tảng chuỗi thời gian đa nhiệm gần đây, ví dụ UniTS \cite{gao2024units},
được huấn luyện trước trên nhiều bộ dữ liệu và nhiều tác vụ (dự báo, phân loại,
imputation), cho phép fine-tune trên một bài toán cụ thể với lượng dữ liệu huấn luyện ít
hơn.

\textbf{Hệ quả cho đồ án:} SAITS được đưa vào như đại diện cho hướng attention-based
imputation, và UniTS (nếu đủ thời gian/tài nguyên) được dùng để kiểm tra liệu fine-tuning
từ một mô hình nền tảng có mang lại lợi thế so với huấn luyện từ đầu trên riêng dữ liệu
Szeged.

\subsection{Đánh giá trên Szeged và các benchmark liên quan}
\label{sec:lit-benchmark-szeged}

\TODO{Rà soát các báo cáo/đồ án/kernel công khai đã dùng bộ dữ liệu Szeged (chủ yếu cho
hồi quy nhiệt độ hoặc phân loại Rain/Snow ở điều kiện dữ liệu đầy đủ), và nêu rõ điểm
khác biệt: các công trình này không giả lập hoặc đánh giá dưới điều kiện mất dữ liệu.}

\subsection{Bảng tổng hợp và khoảng trống nghiên cứu}
\label{sec:lit-tong-hop}

\begin{table}[H]
\centering
\caption{Tổng hợp các hướng tiếp cận liên quan và vị trí của đồ án.}
\label{tab:lit-summary}
\begin{tabular}{@{}p{3.1cm}p{4.3cm}p{4.3cm}@{}}
\toprule
\textbf{Hướng tiếp cận} & \textbf{Giả định chính} & \textbf{Khoảng trống so với đồ án} \\
\midrule
Ngưỡng nhiệt độ \cite{jennings2018,webb2025} & Dữ liệu đầy đủ, mối quan hệ nhiệt độ--pha giáng thủy ổn định & Không xử lý dữ liệu thiếu; độ chính xác thấp trong vùng chồng lấn \\
Đặc trưng + cây quyết định \cite{christ2018tsfresh,lubba2019catch22,chen2016xgboost} & Cửa sổ đầu vào đầy đủ hoặc thiếu ngẫu nhiên nhẹ & Thiếu đánh giá hệ thống theo nhiều tỷ lệ/kiểu thiếu \\
ROCKET \cite{dempster2021minirocket,tan2022multirocket} & Chuỗi đầu vào liên tục, không thiếu & Cần bước nội suy riêng; chưa có benchmark dưới điều kiện thiếu \\
RNN cho dữ liệu thiếu \cite{che2018grud,cao2018brits} & Thiếu dữ liệu là điều kiện vận hành bình thường & Ít được so sánh trực tiếp với ROCKET/đặc trưng cổ điển trên cùng bài toán \\
Attention/mô hình nền tảng \cite{du2023saits,gao2024units} & Có đủ dữ liệu/tài nguyên để huấn luyện hoặc fine-tune & Chi phí tính toán cao, chưa rõ có cần thiết cho bài toán quy mô một trạm \\
\bottomrule
\end{tabular}
\end{table}

Khoảng trống chung là: chưa có nghiên cứu so sánh trực tiếp ba hướng tiếp cận trên cùng
một bộ dữ liệu, cùng một giao thức giả lập thiếu dữ liệu và cùng một tập chỉ số đánh giá
-- đây chính là mục tiêu của đồ án.
